% GENERATED FILE. Edit canonical lessons, then run npm run generate:books.
% canonical-source-hash: fnv1a64:1c5cfbbce36081e0
% canonical-lessons: PA-C254-faram, PA-C254-arzi, PA-C254-sahi, PA-C254-ghalat, PA-C254-mitauna

\chapter{A Form at the Counter}
\label{ch:pa-a2-a-form-at-the-counter}
\hlchaptermodality{254}

\begin{chapteropening}
\textbf{By the end of this chapter:} \emph{I can answer the lines of a simple form.}

Complete the last lesson of chapter 254: I can answer the lines of a simple form.
\end{chapteropening}

\section[\texorpdfstring{\pa{ਫ਼ਾਰਮ} (fāram)}{ਫ਼ਾਰਮ (fāram)}]{\pa{ਫ਼ਾਰਮ} (fāram) — a form to fill in}

\label{lesson:PA-C254-faram}



\begin{quote}
Before the new one: say the Punjabi for to move out of the way, then the Punjabi for an advertisement.
\end{quote}

\subsection*{You'll want to know: \pa{ਫ਼ਾਰਮ}}
\textbf{\pa{ਫ਼ਾਰਮ}} — \emph{fāram} — \textquotedblleft{}a form to fill in\textquotedblright{}.

\begin{quote}
\pa{ਬੈਂਕ} \pa{ਵਿੱਚ} \pa{ਇਹ} \pa{ਫ਼ਾਰਮ} \pa{ਭਰੋ।}
\end{quote}

\begin{quote}
\emph{baiṅk vich eh fāram bharo.}
\end{quote}

\begin{quote}
Fill in this form at the bank.
\end{quote}

A masculine noun.

\begin{cousinweb}
\textbf{\pa{ਫ਼ਾਰਮ}} (\emph{fāram}) is borrowed from English \emph{form}.
\end{cousinweb}

\begin{grammarlens}[title={What you've built}]
One more word for reading notices and messages.
\end{grammarlens}

\subsection*{Your turn}
\begin{itemize}
\raggedright
  \item \emph{Say it:} \emph{fāram}
  \item \emph{Say it:} \emph{fāram}, once more
  \item \emph{Say it:} \emph{baiṅk vich eh fāram bharo.} — Fill in this form at the bank.
  \item \emph{Recall:} say the Punjabi for to move out of the way, then the Punjabi for an advertisement, then say \emph{fāram} again
\end{itemize}

\begin{tcolorbox}[breakable,colback=teal!4,colframe=teal!35!black,title={Before you move on}]
What does \pa{ਫ਼ਾਰਮ} mean? (\textquotedblleft{}A form to fill in\textquotedblright{}.) Say it once more.
\end{tcolorbox}
\section[\texorpdfstring{\pa{ਅਰਜ਼ੀ} (arzī)}{ਅਰਜ਼ੀ (arzī)}]{\pa{ਅਰਜ਼ੀ} (arzī) — an application}

\label{lesson:PA-C254-arzi}



\begin{quote}
Before the new one: say the Punjabi for an advertisement, then the Punjabi for a form to fill in.
\end{quote}

\subsection*{You'll want to know: \pa{ਅਰਜ਼ੀ}}
\textbf{\pa{ਅਰਜ਼ੀ}} — \emph{arzī} — \textquotedblleft{}an application\textquotedblright{}.

\begin{quote}
\pa{ਅਮਨ} \pa{ਨੌਕਰੀ} \pa{ਲਈ} \pa{ਅਰਜ਼ੀ} \pa{ਭੇਜੇਗਾ।}
\end{quote}

\begin{quote}
\emph{Aman naukrī laī arzī bhejegā.}
\end{quote}

\begin{quote}
Aman will send an application for the job.
\end{quote}

A feminine noun.

\begin{cousinweb}
\textbf{\pa{ਅਰਜ਼ੀ}} (\emph{arzī}) is a loan, usually traced to Persian \emph{arzī}, \textquotedblleft{}a petition\textquotedblright{}, itself from Arabic.
\end{cousinweb}

\begin{grammarlens}[title={What you've built}]
One more word for reading notices and messages.
\end{grammarlens}

\subsection*{Your turn}
\begin{itemize}
\raggedright
  \item \emph{Say it:} \emph{arzī}
  \item \emph{Say it:} \emph{arzī}, once more
  \item \emph{Say it:} \emph{Aman naukrī laī arzī bhejegā.} — Aman will send an application for the job.
  \item \emph{Recall:} say the Punjabi for an advertisement, then the Punjabi for a form to fill in, then say \emph{arzī} again
  \item \emph{Write it:} copy \textbf{\pa{ਅਰਜ਼ੀ}} once, checking the whole word
\end{itemize}

\begin{tcolorbox}[breakable,colback=teal!4,colframe=teal!35!black,title={Before you move on}]
What does \pa{ਅਰਜ਼ੀ} mean? (\textquotedblleft{}An application\textquotedblright{}.) Say it once more.
\end{tcolorbox}
\section[\texorpdfstring{\pa{ਸਹੀ} (sahī)}{ਸਹੀ (sahī)}]{\pa{ਸਹੀ} (sahī) — correct}

\label{lesson:PA-C254-sahi}



\begin{quote}
Before the new one: say the Punjabi for a form to fill in, then the Punjabi for an application.
\end{quote}

\subsection*{You'll want to know: \pa{ਸਹੀ}}
\textbf{\pa{ਸਹੀ}} — \emph{sahī} — \textquotedblleft{}correct\textquotedblright{}.

\begin{quote}
\pa{ਕੀ} \pa{ਇਹ} \pa{ਤਾਰੀਖ਼} \pa{ਸਹੀ} \pa{ਹੈ}?
\end{quote}

\begin{quote}
\emph{kī eh tārīkh sahī hai?}
\end{quote}

\begin{quote}
Is this date correct?
\end{quote}

\begin{cousinweb}
\textbf{\pa{ਸਹੀ}} (\emph{sahī}) is a loan, usually traced through Persian to Arabic \emph{sahīh}, \textquotedblleft{}sound, right\textquotedblright{}.
\end{cousinweb}

\begin{grammarlens}[title={What you've built}]
One more word for reading notices and messages.
\end{grammarlens}

\subsection*{Your turn}
\begin{itemize}
\raggedright
  \item \emph{Say it:} \emph{sahī}
  \item \emph{Say it:} \emph{sahī}, once more
  \item \emph{Say it:} \emph{kī eh tārīkh sahī hai?} — Is this date correct?
  \item \emph{Recall:} say the Punjabi for a form to fill in, then the Punjabi for an application, then say \emph{sahī} again
  \item \emph{Write it:} copy \textbf{\pa{ਸਹੀ}} once, checking the whole word
\end{itemize}

\begin{tcolorbox}[breakable,colback=teal!4,colframe=teal!35!black,title={Before you move on}]
What does \pa{ਸਹੀ} mean? (\textquotedblleft{}Correct\textquotedblright{}.) Say it once more.
\end{tcolorbox}
\section[\texorpdfstring{\pa{ਗ਼ਲਤ} (ghalat)}{ਗ਼ਲਤ (ghalat)}]{\pa{ਗ਼ਲਤ} (ghalat) — wrong; mistaken}

\label{lesson:PA-C254-ghalat}



\begin{quote}
Before the new one: say the Punjabi for an application, then the Punjabi for correct.
\end{quote}

\subsection*{You'll want to know: \pa{ਗ਼ਲਤ}}
\textbf{\pa{ਗ਼ਲਤ}} — \emph{ghalat} — \textquotedblleft{}wrong; mistaken\textquotedblright{}.

\begin{quote}
\pa{ਫ਼ਾਰਮ} \pa{ਉੱਤੇ} \pa{ਮੇਰੀ} \pa{ਉਮਰ} \pa{ਗ਼ਲਤ} \pa{ਹੈ।}
\end{quote}

\begin{quote}
\emph{fāram utte merī umar ghalat hai.}
\end{quote}

\begin{quote}
My age is wrong on the form.
\end{quote}

\begin{cousinweb}
\textbf{\pa{ਗ਼ਲਤ}} (\emph{ghalat}) is a loan, usually traced through Persian to Arabic \emph{ghalat}, \textquotedblleft{}an error\textquotedblright{}.
\end{cousinweb}

\begin{grammarlens}[title={What you've built}]
One more word for reading notices and messages.
\end{grammarlens}

\subsection*{Your turn}
\begin{itemize}
\raggedright
  \item \emph{Say it:} \emph{ghalat}
  \item \emph{Say it:} \emph{ghalat}, once more
  \item \emph{Say it:} \emph{fāram utte merī umar ghalat hai.} — My age is wrong on the form.
  \item \emph{Recall:} say the Punjabi for an application, then the Punjabi for correct, then say \emph{ghalat} again
  \item \emph{Write it:} copy \textbf{\pa{ਗ਼ਲਤ}} once, checking the whole word
\end{itemize}

\begin{tcolorbox}[breakable,colback=teal!4,colframe=teal!35!black,title={Before you move on}]
What does \pa{ਗ਼ਲਤ} mean? (\textquotedblleft{}Wrong; mistaken\textquotedblright{}.) Say it once more.
\end{tcolorbox}
\section[\texorpdfstring{\pa{ਮਿਟਾਉਣਾ} (miṭāuṇā)}{ਮਿਟਾਉਣਾ (miṭāuṇā)}]{\pa{ਮਿਟਾਉਣਾ} (miṭāuṇā) — to rub out; to erase}

\label{lesson:PA-C254-mitauna}



\begin{quote}
Before the new one: say the Punjabi for correct, then the Punjabi for wrong.
\end{quote}

\subsection*{You'll want to know: \pa{ਮਿਟਾਉਣਾ}}
\textbf{\pa{ਮਿਟਾਉਣਾ}} — \emph{miṭāuṇā} — \textquotedblleft{}to rub out; to erase\textquotedblright{}.

\begin{quote}
\pa{ਮੈਂ} \pa{ਗ਼ਲਤ} \pa{ਸ਼ਬਦ} \pa{ਮਿਟਾਉਂਦਾ} \pa{ਹਾਂ।}
\end{quote}

\begin{quote}
\emph{maiṁ ghalat shabad miṭāundā hāṁ.}
\end{quote}

\begin{quote}
I rub out the wrong word.
\end{quote}

Past (I …): \pa{ਮੈਂ} \pa{ਗ਼ਲਤ} \pa{ਸ਼ਬਦ} \pa{ਮਿਟਾਇਆ}, \emph{maiṁ ghalat shabad miṭāiā}.

\begin{cousinweb}
\textbf{\pa{ਮਿਟਾਉਣਾ}} (\emph{miṭāuṇā}) is the causative partner of \emph{miṭnā}, \textquotedblleft{}is wiped out\textquotedblright{}, which goes back to Sanskrit \emph{mṛṣṭa}, \textquotedblleft{}rubbed, washed\textquotedblright{}; CDIAL lists Punjabi \emph{miṭnā} under it (CDIAL 10299).
\end{cousinweb}

\begin{grammarlens}[title={What you've built}]
One more word for reading notices and messages.
\end{grammarlens}

\subsection*{Your turn}
\begin{itemize}
\raggedright
  \item \emph{Say it:} \emph{miṭāuṇā}
  \item \emph{Say it:} \emph{miṭāuṇā}, once more
  \item \emph{Say it:} \emph{maiṁ ghalat shabad miṭāundā hāṁ.} — I rub out the wrong word.
  \item \emph{Recall:} say the Punjabi for correct, then the Punjabi for wrong, then say \emph{miṭāuṇā} again
  \item \emph{Write it:} copy \textbf{\pa{ਮਿਟਾਉਣਾ}} once, checking the whole word
\end{itemize}

\begin{tcolorbox}[breakable,colback=teal!4,colframe=teal!35!black,title={Before you move on}]
What does \pa{ਮਿਟਾਉਣਾ} mean? (\textquotedblleft{}To rub out; to erase\textquotedblright{}.) Say it once more.
\end{tcolorbox}
